
二十世纪六十年代，V.I.~\trArnold{}在\textbf{\trSymplecticTopology}%
（symplectic topology）中宣布了若干影响深远且硕果累累的猜想。本章致力于%
他关于\trSymplectic\trMapping{}之\trFixedPoint{}的那个猜想，它起源于一个%
要追溯到世纪之交的\trCelestialMechanics{}问题。用\trDynamics{}的方式重新表述%
之后，这个猜想以\trCompact\trSymplecticManifold{}的\trTopology{}来给出其%
上一个\trTimeDependent\trHamiltonianVectorField{}的整体%
\textbf{\trForcedOscillation}（forced oscillation）数目的下界。%
\trForcedOscillation{}由\trAction{}原理挑出；有关这个猜想的求解则在%
\trCalculusOfVariations{}和\trNonlinear\trElliptic{}系统中催生了新的思想%
与技巧。特别地，%
它促使 A.~\trFloer{}于 1986 年构造出他的\trHomology{}理论。

下文中我们将描述\trArnold{}猜想的历史，并在 6.2 节中从标准\trTorus{}%
$T^{2n}$ 这一特殊情形的证明开始。这里我们可以使用第 3 章中引入的 $H^{1/2}$%
框架下的\trAction{}原理，但这一次的目标是求出它的全部\trCriticalPoint。%
我们的策略受到 C.~\trConley{}研究\trDynamicalSystem{}的\trTopology{}进路的%
启发：我们将研究\textbf{\trRegularization\trGradient{}方程}（regularized %
gradient equation）之全体\trBounded{}解所成集合的结构，该方程正是由定义在%
$T^{2n}$ 的诸\trContractible\trLoop{}之上的那个\trActionFunctional{}所产生的。这个%
\trInvariant{}集由诸\trCriticalPoint{}连同它们的全部%
\textbf{\trConnectingOrbit}（connecting orbit）组成。这样一来，无穷维%
\trLoopSpace{}中\trGradient{}流的研究便被化归为一个\trCompact\trMetric{}%
空间的\trGradientLike{}连续\trFlow{}的研究，后者的\textbf{\trRestPoint}%
（rest point）正是所需的那些\trCriticalPoint。事实表明，这个\trCompact{}%
空间含有该\trTorus{}的\textbf{\trCohomology}（cohomology）。\trRestPoint%
的数目随后便可由 6.3 节中介绍的\trLjusternik-\trSchnirelman{}理论来估计。

把\trTorus{}情形之证明背后的几何思想用\trPartialDifferentialEquation{}的%
语言%
重新解释，将在 6.4 节中把我们引向\trArnold{}猜想对更大一类满足%
$[\omega]\big|_{\pi_{2}(M)}=0$ 的\trCompact\trSymplecticManifold{}%
$(M,\omega)$ 的证明。在这个一般情形中没有自然的\trRegularization，我们%
不得不去研究未\trRegularization{}的\trGradient{}系统之\trBounded{}解的%
结构。它们是一类特殊的\trCauchyRiemann{}型一阶\trElliptic%
\trPartialDifferentialEquation{}组的\trSmooth{}解。这些解与%
M.~\trGromov{}在 $(M,\omega)$ 中的\trPseudoholomorphic\trCurve{}有关。%
解集合的至关紧要的\trCompactness{}将借助一种分析技巧来建立，它有时被称为%
\textbf{\trBubblingOffAnalysis}（bubbling off analysis）。为了洞悉解集合%
的\trTopology，我们将使用一种\trApproximation{}步骤：我们用某个特殊%
\trElliptic\trBoundaryValueProblem{}的解来\trApproximation{}这些%
\trBounded{}解，对这个%
边值问题可以应用基于\trCauchyRiemann{}算子的\trFredholm{}理论。利用%
\trAlexander-\trSpanier{}上同调的\trContinuity，我们随后便如%
\trTorus{}情形那样，%
借\trLjusternik-\trSchnirelman{}理论推出\trArnold{}猜想。

此外，若假设一切\trForcedOscillation{}都是\trNondegenerate{}的，这些几何%
思想将在 6.5 节中把我们引向 A.~\trFloer{}对\trMorseTheory{}与\trFloer{}%
同调的全新进路。它使用了未\trRegularization{}的\trGradient{}方程之%
\trBounded{}解集合的一个附加结构：即，由于 $M$ 是\trCompact{}的，%
\trNondegenerate{}的\trForcedOscillation{}只有有限多个，而且\trGradient{}%
方程的每个\trBounded{}解在正向与反向时间上都渐近（指数式）地趋于该%
\trHamiltonianSystem{}的一个特定的\trForcedOscillation。这些一维的%
\trConnectingOrbit{}确定了一个由诸\trForcedOscillation{}生成的代数复形，%
其分次由一个与%
\trNondegenerate\trForcedOscillation{}相联系的\trMaslov{}型指数定义，该指数%
扮演着\trMorseIndex{}的角色。对\trFloer{}美妙的思想，我们仅勾画其轮廓。%
把\trFloer{}的构造与我们第 3 章中特殊\trCapacity{} $c_{0}$ 的构造相结合，%
便产生一种\trSymplectic\trHomology{}理论，然而它尚未定型。我们将在 6.6 节%
中对它不加证明地勾画。这些理论在技术上的要求——特别是繁难的%
\textbf{\trTransversality}（transversality）论证——相当高深，超出了本书%
的范围。与前面各章相反，第 6 章并不是自足的。
